\documentclass[10pt,a4paper]{amsart}
\usepackage[margin=19mm]{geometry}
\usepackage{amsmath,amssymb,mathtools}
\usepackage[scr=rsfs,cal=euler]{mathalfa}
\usepackage{xfrac}
\usepackage[unicode,hidelinks,bookmarks=false]{hyperref}
\newcommand{\ie}{i.e.\ }
\newcommand{\cf}{cf.\ }
\newcommand{\resp}{respectively\ }
\newcommand{\Subsection}{\subsection}
\providecommand{\nobreakdash}{\nobreak\hbox{-}}
\setlength{\emergencystretch}{2em}
\multlinegap0pt
\usepackage{amssymb}
\usepackage{mathtools}
\usepackage{algorithm}\usepackage{algpseudocode}
\newtheorem{theorem}{Theorem}
\newcommand{\E}{\mathbb{E}}
\newcommand{\bV}{\mathbf{V}}
\newcommand{\bx}{\mathbf{x}}
\newcommand{\by}{\mathbf{y}}
\newcommand{\kbar}{\bar{k}}
\title{Drifting to Boltzmann:\\ Million-Fold Acceleration in Boltzmann Sampling with Force-Guided Drifting}
\author{Pipi Hu}
\date{Source: arXiv:2603.05527v1 (CC BY 4.0)}
\begin{document}
\maketitle

\noindent\emph{Source and licence.} Drifting to Boltzmann:\\ Million-Fold Acceleration in Boltzmann Sampling with Force-Guided Drifting. Version arXiv:2603.05527v1 (CC BY 4.0), distributed by arXiv under the Creative Commons Attribution 4.0 International licence, http://creativecommons.org/licenses/by/4.0/, as recorded in the arXiv licence metadata for this exact version and shown on the abstract page https://arxiv.org/abs/2603.05527v1. Full licence notice: \texttt{licenses/arxiv-2603.05527-CC-BY-4.0.txt}.\par\smallskip

\noindent\textit{Editorial scope.} Counted unit: the article's theorem thm:main (source0.tex lines 213--223). The article's complete original proof of it is delivered with it (source0.tex lines 225--251, one \texttt{proof} environment of 27 lines). No mathematics is written, repaired or re-derived: the statement and the proof are the article's own lines, in the article's own notation, and no retained line is substituted or re-flowed.  No further source-local context is needed: every cross-reference of every retained line resolves inside the two delivered spans.  Nothing was omitted from any retained span, so the package carries no omission record.  No retained line contains a citation, so the wrapper carries no bibliography and no bibliography entry.  No retained line contains a figure environment, an image-inclusion command or drawing code, so no image is delivered and the package declares no asset dependency.  Editorial formatting only, in addition: an AMS \texttt{amsart} preamble that restates the article's own declarations and macros used by the retained lines, namely the environments \texttt{theorem} on separate counters, together with the standard packages those lines need.  The retained environments print in this document as Theorem 1 in order of appearance, whereas the article prints the same items with the numbering of its own full document.  The complete native source of the article as distributed in the arXiv e-print for this exact version is bundled unchanged as \texttt{sources/195808/source0.tex}.
\begin{theorem}[Drifting Score Identity]
\label{thm:main}
Let $k(\bx,\by) = \exp(-\|\bx-\by\|^2/(2\tau^2))$ be a Gaussian kernel.
For any sufficiently regular distribution $p$ (i.e., $p(\by)k(\bx,\by) \to 0$ as $\|\by\|\to\infty$), the attraction term of the drifting field satisfies, in the population limit:
\begin{equation}
    \bV_p^+(\bx) = \tau^2 \cdot \E_p\!\left[\kbar(\bx,\by)\,\nabla_\by \log p(\by)\right],
    \label{eq:main-identity}
\end{equation}
where $\kbar(\bx,\by) = k(\bx,\by)/\E_p[k(\bx,\by')]$ is the normalized kernel weight.
That is, the attraction term equals $\tau^2$ times a kernel-weighted average of the score function $\nabla \log p$ of the sampling distribution.
\end{theorem}

\begin{proof}
The attraction term in the population limit is
$\bV_p^+(\bx) = {\E_{p(\by)}[k(\bx,\by)(\by - \bx)]}/{\E_{p(\by)}[k(\bx,\by)]}$.
We transform the numerator.

\textit{Step 1: Kernel gradient identity.}
For the Gaussian kernel, $\nabla_\by k(\bx,\by) = -k(\bx,\by)(\by - \bx)/\tau^2$, so:
\begin{equation}
    k(\bx,\by)(\by - \bx) = -\tau^2 \nabla_\by k(\bx,\by).
    \label{eq:kernel-grad}
\end{equation}

\textit{Step 2: Integration by parts.}
Substituting~\eqref{eq:kernel-grad} into the numerator and applying integration by parts (boundary terms vanish by the regularity assumption):
\begin{align}
    \E_p[k(\bx,\by)(\by - \bx)]
    &= -\tau^2\int p(\by)\,\nabla_\by k(\bx,\by)\,d\by \nonumber \\
    &= \tau^2\int \nabla_\by p(\by) \cdot k(\bx,\by)\,d\by
    = \tau^2\,\E_p\!\left[\nabla_\by \log p(\by) \cdot k(\bx,\by)\right].
    \label{eq:ibp}
\end{align}

Dividing by $\E_p[k(\bx,\by)]$:
\begin{equation}
    \bV_p^+(\bx) = \tau^2 \cdot \frac{\E_p[\nabla_\by \log p(\by) \cdot k(\bx,\by)]}{\E_p[k(\bx,\by)]} = \tau^2 \cdot \E_p\!\left[\kbar(\bx,\by)\,\nabla_\by \log p(\by)\right]. \qedhere
\end{equation}
\end{proof}

\end{document}
